\documentclass[11pt,a4paper]{article}
\usepackage[utf8]{inputenc}
\usepackage[T1]{fontenc}
\usepackage{lmodern}
\usepackage[french]{babel}
\usepackage{amsmath,amssymb,amsthm}
\usepackage[margin=2.2cm,top=2cm,bottom=2.2cm]{geometry}
\usepackage{enumitem}
\usepackage{fancyhdr}
\usepackage{array}
\usepackage[hidelinks]{hyperref}
\setlist[enumerate,1]{label=\arabic*.,leftmargin=*,itemsep=0.6em}
\setlist[enumerate,2]{label=(\alph*),leftmargin=1.8em,itemsep=0.3em}
\newcommand{\R}{\mathbb{R}}
\newcommand{\N}{\mathbb{N}}
\newcommand{\C}{\mathbb{C}}
\newcommand{\sh}{\operatorname{sh}}
\newcommand{\ch}{\operatorname{ch}}
\newcommand{\ee}{\mathrm{e}}
\newcommand{\dd}{\,\mathrm{d}}
\pagestyle{fancy}\fancyhf{}
\fancyhead[L]{Grille de rédaction}
\fancyhead[R]{sans indice, sans corrigé}
\fancyfoot[C]{\thepage}
\renewcommand{\headrulewidth}{0.3pt}
\newcommand{\exo}[3]{\vspace{1.1em}\noindent\textbf{\large #1}\hfill\textit{#2 \,·\, #3}\par\smallskip}
\begin{document}
\begin{center}
{\LARGE\bfseries Grille de rédaction}\\[0.3em]
{\large à garder à côté de ta copie, à relire avant de rendre chaque question}
\end{center}

\noindent\fbox{\parbox{\dimexpr\linewidth-2\fboxsep-2\fboxrule\relax}{\small Ces 10 points viennent de \textbf{tes} copies de la feuille 1 : ce sont tes erreurs, pas des généralités.}}


\begin{enumerate}[label=\textbf{\arabic*.},itemsep=0.9em]
\item \textbf{Négation et quantificateurs.} Avant une preuve par l'absurde, écris la négation \emph{exacte}. « Pas de même signe à partir d'un rang » ne veut pas dire « toujours de signe opposé ».
\item \textbf{Bornes d'une intégrale.} Une relation de Chasles n'est permise que si l'intégrande est défini et continu sur tout l'intervalle utilisé (ex. : $1/\ln t$ n'existe pas en $t=1$).
\item \textbf{Dérivée d'une composée.} Écris la formule ($F(x^2)'=2x\,F'(x^2)$) puis \emph{simplifie} ($\ln(x^2)=2\ln x$). Vérifie sur une valeur numérique.
\item \textbf{« $\sim$ » n'est pas une égalité.} On ne soustrait ni n'additionne des équivalents, on ne remplace pas un équivalent dans une différence. Pour une différence : un DL.
\item \textbf{Ordre d'un DL.} Si on multiplie par $n$ (ou $x$), garde un terme de plus. Ne tronque jamais une expression intermédiaire avant d'avoir regardé par quoi elle sera multipliée.
\item \textbf{Coefficients usuels.} Récite-les avant de les utiliser ($\sqrt{1+u}=1+\tfrac u2-\tfrac{u^2}8+\cdots$) : relis chaque DL une fois écrit.
\item \textbf{Limites écrites avec leur flèche.} Pas de « $1=n(u_n+1)$ en $+\infty$ » : écris $n(u_n+1)\longrightarrow 1$, donc $u_n+1\sim\frac1n$.
\item \textbf{« À partir d'un certain rang ».} Une inégalité vraie seulement pour $n\ge 3$ doit être dite ainsi. Et pour un critère d'équivalence sur des séries : vérifie la \emph{positivité} (ou passe aux modules).
\item \textbf{Croissances comparées précises.} $\ln n=o(n)$ ne suffit pas pour $\sum\frac{\ln n}{n^{3/2}}$ : il faut $\ln n=o(n^{\varepsilon})$ avec un $\varepsilon$ bien choisi.
\item \textbf{Contrôle de vraisemblance (10 secondes).} Signe, ordre de grandeur, $n=2$ ou $x=2$. Et relis chaque recopie : $t-1$ n'est pas $t+1$.
\end{enumerate}

\vfill
\begin{small}
\noindent\textbf{Sources.} Établie à partir des copies manuscrites de la feuille 1 (21 octobre 2026) et vérifiée par calcul numérique.
\end{small}
\end{document}
